\section{Event Study: Home Sale Prices}
\label{sec:empirical_results_home_sales}

\subsection{Design}

If the supply restrictions documented in the previous two sections are economically meaningful, they should eventually be capitalized into housing prices. To test that idea, I apply the same redistricting variation from Section~\ref{sec:empirical_results_permits} to home sale prices.

The specification is estimated at the transaction level:
%
\begin{equation}
\label{eq:event_study_sales}
\ln(\text{Price}_{iy}) = \alpha_{gs} + \lambda_{gy} + X_{iy}'\gamma + \sum_{\tau \neq -1} \beta_\tau \left( \Delta\text{Stringency}_i \times \mathbf{1}[\text{event time} = \tau] \right) + \varepsilon_{iy}
\end{equation}
%
where $\ln(\text{Price}_{iy})$ is the log sale price for transaction $i$ in sale year $y$, $\alpha_{gs}$ is a border-pair-by-origin-side fixed effect, $\lambda_{gy}$ is a border-pair-by-year fixed effect, and $X_{iy}$ includes hedonic controls for building area, lot area, building age, bedrooms, bathrooms, garage presence, and parcel-level distances to the nearest school, park, major road, and Lake Michigan. The sample is restricted to transactions within 1000ft of cohort-specific pre-redistricting ward boundaries with uniform weights. As in the permit panel, switched blocks use the fixed origin-destination boundary created by redistricting, while non-switched controls use their nearest incident pre-redistricting boundary. Standard errors are clustered at the census block level. I focus on the 2015 redistricting cohort because the 2023 redistricting provides too little post-period data for meaningful inference. The sample window runs from 2010 to 2020, providing roughly five years on each side. Appendix Table~\ref{tab:redistricting_sample_support} reports the relevant block counts and regression sample size.

The fixed-effects structure here differs from the permit specification in Section~\ref{sec:empirical_results_permits}, and the reason lies in the outcome data. Permits are counts that aggregate naturally to a balanced block-year panel: every block is observed in every year, even if the count is zero. Home sales, by contrast, are individual transactions with property-level hedonic characteristics such as building area, age, and lot size that vary across properties within the same block and matter for price. Many blocks have zero or one sale in a given year, so collapsing to block-year level would throw away the hedonic variation that helps isolate the price effect from compositional changes. I therefore use border-pair-by-side fixed effects, which group blocks by their pre-redistricting ward and boundary location, together with border-pair-by-year fixed effects that absorb local time-varying shocks. Identification thus comes from comparing treated and control blocks on the same side of the same boundary in the same year, rather than from within-block variation over time as in the permit specification. I keep the same 1000ft window as in the permit design. The parallel-trends assumption is correspondingly stronger: treated and control blocks must be on similar price trajectories absent redistricting. The pre-period coefficients below do not show a clear directional trend, though the evidence is noisier than in the permit design.
Appendix~\ref{sec:appendix_event_study} reports additional rent and home-sales border-RD checks using the updated amenity controls.

\subsection{Results}

Figure~\ref{fig:event_study_sales} presents the event study with the full hedonic-and-amenity control set. The pre-period coefficients move around somewhat but are generally centered near zero. Following redistricting, the post-period coefficients remain imprecise and do not show a clear price response on blocks assigned to more stringent aldermen.

Figure~\ref{fig:event_study_sales_split} shows the continuous split, allowing separate slopes for blocks that experienced positive versus negative changes in stringency. The split estimates are also imprecise and do not show a stable symmetric pattern between blocks moving to more stringent aldermen and blocks moving to more lenient ones.

\begin{figure}[htbp]
    \centering
    \includegraphics[width=0.8\textwidth]{../tasks/run_event_study_sales_disaggregate/output/event_study_disaggregate_yearly_cohort_2015_continuous_uniform_1000ft_amenity_full_geo_wardpair_clust_block.pdf}
    \caption{Event Study: Effect of Alderman Stringency on Home Sale Prices}
    \figurenotes{Continuous treatment: effect of 1 SD increase in stringency. 2015 redistricting cohort. Border-pair-by-side and border-pair-by-year fixed effects. Uniform weights with 1000ft bandwidth. Includes hedonic controls for building area, lot area, building age, bedrooms, bathrooms, and garage, plus parcel-level distances to the nearest school, park, major road, and Lake Michigan. Year $-1$ omitted as reference. 95\% confidence intervals based on standard errors clustered at the census block level. The pooled post-period estimate in the full-control specification is $-0.001$ (SE $= 0.009$), and the joint pre-trend test gives $p = 0.454$.}
    \label{fig:event_study_sales}
\end{figure}

\begin{figure}[htbp]
    \centering
    \includegraphics[width=0.8\textwidth]{../tasks/run_event_study_sales_disaggregate/output/event_study_disaggregate_yearly_cohort_2015_continuous_split_uniform_1000ft_amenity_full_geo_wardpair_clust_block.pdf}
    \caption{Home Sale Event Study: Continuous Split by Treatment Direction}
    \figurenotes{Same specification as Figure~\ref{fig:event_study_sales}, estimated separately for blocks with $\Delta\text{Stringency} > 0$ (moved to more stringent) and $\Delta\text{Stringency} < 0$ (moved to more lenient). 95\% confidence intervals based on standard errors clustered at the census block level.}
    \label{fig:event_study_sales_split}
\end{figure}

Table~\ref{tab:did_sales_2015_current} collapses the post-period into a single coefficient. The estimates are close to zero across specifications: $-0.3\%$ without controls, $0.03\%$ with hedonic controls, and $-0.05\%$ after adding parcel-level amenity distance controls. The standard errors are about 0.8--1.0 percentage points, so the transaction-level design does not provide clear evidence of capitalization over this window.

\begin{table}[htbp]
    \caption{Effect of Alderman Stringency on Home Sale Prices}
    \label{tab:did_sales_2015_current}

    \input{../tasks/run_event_study_sales_disaggregate/output/did_table_sales_2015_uniform_1000ft_geo_wardpair_clust_block.tex}

    \footnotesize
    \textit{Notes:} Transaction-level regressions of log sale price using the 2015 implementation cohort, restricted to transactions within 1000ft of ward boundaries and relative years $-5$ through $5$. All columns use the same transaction sample. Column (2) adds hedonic controls for building area, lot area, building age, bedrooms, bathrooms, and garage. Column (3) further adds parcel-level distances to the nearest school, park, major road, and Lake Michigan. Individual control coefficients are omitted for readability. Uniform weighting. Standard errors clustered at the block level. * $p<0.10$, ** $p<0.05$, *** $p<0.01$.
\end{table}

\subsection{Interpretation}

The null price result should not be read as evidence that supply restrictions have no welfare cost. The most direct effects documented above occur in permitting and realized density. Capitalization into transaction prices may require a longer horizon, may be muted by limited local turnover, or may be hard to detect in a transaction-level design with changing sale composition.

For that reason, the home-sales evidence plays a different role from the density and permit results. It is a test of whether the detected supply-side effects show up in nearby sales prices over the available event-study window. In the current sample, they do not. The paper's main empirical conclusion therefore rests on the two supply-side findings: stringency reduces density at ward boundaries (Section~\ref{sec:empirical_results_density}) and reduces the flow of high-discretion permits when the alderman changes (Section~\ref{sec:empirical_results_permits}).
